% ============================================================
%  Plantilla para el Proyecto Final de Aprendizaje Automático
%  Escuela Politécnica Superior — Universidad de Alicante
%  Formato: artículo científico (una columna)
% ============================================================

\documentclass[11pt,a4paper]{article}

\usepackage[utf8]{inputenc}
\usepackage[T1]{fontenc}
\usepackage[spanish,es-tabla]{babel}
\usepackage[margin=2.5cm]{geometry}
\usepackage{amsmath}
\usepackage{amssymb}
\usepackage{bm}
\usepackage{graphicx}
\usepackage{booktabs}
\usepackage{multirow}
\usepackage{array}
\usepackage{float}
\usepackage{listings}
\usepackage{xcolor}

\lstset{
    language=Python,
    basicstyle=\ttfamily\footnotesize,
    keywordstyle=\color{blue},
    commentstyle=\color{gray},
    stringstyle=\color{red!70!black},
    showstringspaces=false,
    breaklines=true,
    frame=single,
    numbers=left,
    numberstyle=\tiny\color{gray}
}

\usepackage{hyperref}
\hypersetup{
    colorlinks=true,
    linkcolor=blue!70!black,
    citecolor=blue!70!black,
    urlcolor=blue!70!black
}

% Configuración de bloques de orientación
\usepackage{mdframed}

\newmdenv[
  backgroundcolor=blue!6,
  linecolor=blue!50!black,
  linewidth=1pt,
  roundcorner=4pt,
  innerleftmargin=8pt,
  innerrightmargin=8pt,
  innertopmargin=6pt,
  innerbottommargin=6pt,
  skipabove=6pt,
  skipbelow=4pt
]{orientacion}

\newcommand{\orientaciontitle}{%
  {\small\bfseries\color{blue!60!black}Orientación}\par\smallskip
}

% Otros 
\usepackage{microtype}


% ============================================================
%  DATOS DEL TRABAJO 
% ============================================================

\title{
    \textbf{[Título del trabajo]}\\[4pt]
    \large [Subtítulo opcional]
}

\author{
    \textbf{[Nombre Apellido1]}\thanks{\texttt{[email1@alu.ua.es]}}
    \quad
    \textbf{[Nombre Apellido2]}\thanks{\texttt{[email2@alu.ua.es]}}
    \quad
    \textbf{[Nombre Apellido3]}\thanks{\texttt{[email3@alu.ua.es]}}
    \\[4pt]
    \small Grado en Ingeniería en Inteligencia Artificial
    --- Escuela Politécnica Superior, Universidad de Alicante\\
    \small Aprendizaje Avanzado --- Curso 2025/2026
}

\date{}

% ============================================================

\begin{document}

\maketitle
\thispagestyle{empty}

% Abstract 

\begin{abstract}
\textit{Resumen del trabajo en 150--250 palabras. Debe incluir:
(1)~el problema abordado y su relevancia real,
(2)~el \textit{dataset} utilizado,
(3)~las problemáticas identificadas y las técnicas aplicadas,
(4)~los resultados principales.
\textbf{Escríbelo al final, una vez completado el trabajo.}}
\end{abstract}

\smallskip
\noindent\textbf{Palabras clave:} aprendizaje automático, [problemática 1], [problemática 2], [dataset], [otra palabra clave].

\tableofcontents
\newpage


% ============================================================
\section{Introducción}
% ============================================================

\begin{orientacion}
\orientaciontitle
\textit{La introducción debe responder a cuatro preguntas: (1)~?`Cuál es el problema y por qué es relevante? (2)~?`Qué implicaciones éticas, sociales o legales tiene? (3)~?`Qué se propone hacer este trabajo? (4)~?`Cómo se organiza el resto del documento? No es un resumen del trabajo, sino una motivación y contextualización.}
\end{orientacion}

[Introducción aquí: 3--4 párrafos.]

El resto del artículo se organiza como sigue.
La Sección~\ref{sec:related} revisa el trabajo relacionado.
La Sección~\ref{sec:datos} describe el \textit{dataset} y el análisis exploratorio.
La Sección~\ref{sec:metodologia} detalla la metodología.
La Sección~\ref{sec:resultados} presenta los resultados experimentales.
La Sección~\ref{sec:discusion} discute las implicaciones.
La Sección~\ref{sec:conclusiones} concluye el trabajo.


% ============================================================
\section{Trabajo Relacionado}
\label{sec:related}
% ============================================================

\begin{orientacion}
\orientaciontitle
\textit{No se trata de resumir trabajos uno a uno, sino de sintetizar la literatura de forma crítica: ?`qué se ha hecho antes?, ?`qué limitaciones tienen esos enfoques?, ?`cómo se sitúa este trabajo respecto al estado del arte? Organiza la sección por temática en lugar de por autor. Cita al menos 5--6 referencias de calidad~\cite{ref1}.}
\end{orientacion}

\subsection{[Trabajos sobre el problema abordado]}

[Síntesis de trabajos previos sobre el mismo problema o \textit{dataset}~\cite{ref1, ref2}.]

\subsection{[Técnicas para la problemática 1]}

[Resumen del estado del arte en la primera problemática identificada~\cite{ref3}.]

\subsection{[Técnicas para la problemática 2]}

[Resumen del estado del arte en la segunda problemática identificada~\cite{ref4}.]


% ============================================================
\section{Descripción del Problema y los Datos}
\label{sec:datos}
% ============================================================

\subsection{Definición formal de la tarea}

\begin{orientacion}
\orientaciontitle
\textit{Define formalmente la tarea de aprendizaje: tipo de tarea (clasificación, regresión\ldots), espacio de entrada y salida, y función objetivo. Por ejemplo: dado un vector de \textit{features} $\bm{x} \in \mathbb{R}^d$, aprender $f\colon \mathbb{R}^d \to \mathcal{Y}$ que minimiza $\mathcal{L}(f) = \mathbb{E}_{(\bm{x},y)\sim P}[\ell(f(\bm{x}),y)]$.}
\end{orientacion}

[Definición formal de la tarea aquí.]

\subsection{Dataset}

\begin{orientacion}
\orientaciontitle
\textit{Describe el \textit{dataset}: origen, licencia, número de instancias, número de \textit{features}, variable objetivo y distribución de clases. Incluye también las consideraciones éticas sobre su uso: ?`quiénes son los individuos representados?, ?`cómo se recogieron los datos?, ?`hay sesgos de selección conocidos?}
\end{orientacion}

La Tabla~\ref{tab:dataset} resume las características principales del \textit{dataset} utilizado.

\begin{table}[H]
\centering
\caption{Características del \textit{dataset}.}
\label{tab:dataset}
\small
\begin{tabular}{@{}ll@{}}
\toprule
\textbf{Característica} & \textbf{Valor} \\
\midrule
Fuente / enlace           & [UCI / Kaggle / \ldots] \\
Número de instancias  & \\
Número de \textit{features} & \\
Variable objetivo         & \\
Distribución de clases & [X\,\% clase 0 / Y\,\% clase 1] \\
Valores ausentes          & [Sí / No --- estrategia de imputación] \\
Licencia                  & \\
\bottomrule
\end{tabular}
\end{table}

\subsection{Análisis exploratorio (EDA)}

\begin{orientacion}
\orientaciontitle
\textit{Incluye las visualizaciones más informativas: distribución de la variable objetivo, correlaciones relevantes, distribución de variables sensibles si las hay. El EDA debe estar orientado a identificar las problemáticas del trabajo, no ser un catálogo exhaustivo de gráficas.}
\end{orientacion}

[Análisis exploratorio aquí: 2--3 párrafos con las figuras más relevantes.]

\begin{figure}[H]
\centering
% \includegraphics[width=0.7\textwidth]{figuras/eda.pdf}
\textit{[Insertar figura del EDA.]}
\caption{[Descripción de la figura del EDA.]}
\label{fig:eda}
\end{figure}

\subsection{Problemáticas identificadas}

\begin{orientacion}
\orientaciontitle
\textit{Justifica aquí qué problemáticas son relevantes para este problema concreto. La argumentación debe basarse en: (1)~el contexto de uso real del sistema (?`quién usa el modelo?, ?`sobre quién toma decisiones?), (2)~las características del \textit{dataset} (?`hay atributos sensibles?, ?`está desbalanceado?), y (3)~las obligaciones legales aplicables (RGPD, AI Act).}
\end{orientacion}

A partir del análisis del \textit{dataset} y del contexto de uso del sistema, se han identificado las siguientes problemáticas:

\begin{itemize}
    \item \textbf{[Problemática 1]:} [Justificación.]
    \item \textbf{[Problemática 2]:} [Justificación.]
\end{itemize}


% ============================================================
\section{Metodología}
\label{sec:metodologia}
% ============================================================

\begin{orientacion}
\orientaciontitle
\textit{Esta sección debe ser suficientemente detallada para que el trabajo sea reproducible. Describe el \textit{pipeline} completo, justifica cada decisión de dise\~{n}o, e incluye la formulación matemática de las técnicas que lo requieran.}
\end{orientacion}

\subsection{Preprocesamiento}

[Cadena de preprocesamiento: imputación de valores ausentes, escalado, codificación de variables categóricas, división train/test y estratificación.]

\subsection{Modelo \textit{baseline}}

\begin{orientacion}
\orientaciontitle
\textit{El \textit{baseline} es el punto de referencia contra el que compararás las técnicas aplicadas. Elige un modelo estándar apropiado para la tarea y justifica su elección.}
\end{orientacion}

[Descripción del modelo \textit{baseline} y justificación.]

\subsection{Tratamiento de la problemática 1: [nombre]}

\begin{orientacion}
\orientaciontitle
\textit{Describe la técnica aplicada. Incluye la formulación matemática cuando sea necesaria para entender el método. Justifica por qué esta técnica es adecuada para el problema y el \textit{dataset} elegidos.}
\end{orientacion}

[Descripción, formulación y justificación.]

% Ejemplo de ecuación numerada:
% \begin{equation}
%   \min_{w} \mathcal{L}(w) \quad \text{s.a.} \quad
%   \bigl|P(\hat{Y}=1 \mid A=a) - P(\hat{Y}=1 \mid A=b)\bigr| \leq \epsilon
%   \label{eq:ejemplo}
% \end{equation}

\subsection{Tratamiento de la problemática 2: [nombre]}

[Descripción, formulación y justificación.]

\subsection{Protocolo de evaluación}

\begin{orientacion}
\orientaciontitle
\textit{Describe cómo se evalúan los modelos: tipo de validación cruzada, métricas elegidas y justificación de su elección. Si el \textit{dataset} está desbalanceado, justifica el uso de métricas robustas (F1, AUC-ROC, etc.) en lugar de \textit{accuracy}.}
\end{orientacion}

[Descripción del protocolo de evaluación y justificación de las métricas elegidas.]


% ============================================================
\section{Experimentos y Resultados}
\label{sec:resultados}
% ============================================================

\begin{orientacion}
\orientaciontitle
\textit{Esta sección presenta los resultados sin interpretarlos (la interpretación va en la Sección~\ref{sec:discusion}). Usa tablas comparativas que permitan ver claramente la diferencia entre el \textit{baseline} y cada intervención. Reporta siempre media $\pm$ desviación estándar cuando uses validación cruzada.}
\end{orientacion}

\subsection{Configuración experimental}

[Entorno de experimentación: versiones de librerías principales y semilla aleatoria fijada.]

\begin{lstlisting}[caption={Semilla aleatoria y versiones principales.}]
import numpy as np
SEED = 42
np.random.seed(SEED)
# sklearn.__version__ = X.X.X
\end{lstlisting}

\subsection{Resultados comparativos}

La Tabla~\ref{tab:resultados} compara el rendimiento del \textit{baseline} con los modelos que incorporan las técnicas para cada problemática.

\begin{table}[H]
\centering
\caption{Comparativa de resultados (media $\pm$ desv. estándar, validación cruzada 5~\textit{folds}).}
\label{tab:resultados}
\small
\begin{tabular}{@{}lcccc@{}}
\toprule
\textbf{Modelo / Configuración} & \textbf{AUC-ROC} & \textbf{F1} & \textbf{Métrica específica} & \textbf{Tiempo (s)} \\
\midrule
\textit{Baseline}        & $0.XX \pm 0.XX$ & $0.XX \pm 0.XX$ & --- & \\
+ Problemática 1     & $0.XX \pm 0.XX$ & $0.XX \pm 0.XX$ & $0.XX$ & \\
+ Problemática 2     & $0.XX \pm 0.XX$ & $0.XX \pm 0.XX$ & $0.XX$ & \\
\bottomrule
\end{tabular}
\end{table}

\subsection{Resultados específicos por problemática}

[Figuras o tablas específicas de cada problemática: \textit{summary plot} de SHAP, curva privacidad--utilidad, comparativa de métricas de \textit{fairness}, etc.]

\begin{figure}[H]
\centering
% \includegraphics[width=0.7\textwidth]{figuras/resultado_p1.pdf}
\textit{[Insertar figura específica de la problemática 1.]}
\caption{[Descripción del resultado de la problemática 1.]}
\label{fig:resultado_p1}
\end{figure}

\begin{figure}[H]
\centering
% \includegraphics[width=0.7\textwidth]{figuras/resultado_p2.pdf}
\textit{[Insertar figura específica de la problemática 2.]}
\caption{[Descripción del resultado de la problemática 2.]}
\label{fig:resultado_p2}
\end{figure}


% ============================================================
\section{Discusión}
\label{sec:discusion}
% ============================================================

\begin{orientacion}
\orientaciontitle
\textit{Interpreta los resultados en el contexto del problema real: ?`qué significan para los usuarios afectados por el sistema? Analiza las limitaciones con rigor y evita afirmaciones genéricas sobre ética; las implicaciones deben ser específicas y estar fundamentadas en los resultados obtenidos.}
\end{orientacion}

\subsection{Interpretación de los resultados}

[Interpreta los resultados en el contexto del problema real: 2--3 párrafos.]

\subsection{Limitaciones}

[Identifica y analiza las limitaciones no triviales: supuestos asumidos, restricciones del \textit{dataset}, limitaciones de las técnicas aplicadas.]

\subsection{Implicaciones éticas, legales y sociales}

[Reflexión sobre las implicaciones reales del sistema. ?`Cómo afectan los resultados a los grupos implicados? ?`Qué obligaciones legales aplicarían bajo el RGPD o el AI Act?]


% ============================================================
\section{Conclusiones}
\label{sec:conclusiones}
% ============================================================

\begin{orientacion}
\orientaciontitle
\textit{Síntesis de las contribuciones en 3--4 párrafos. Lecciones aprendidas. Líneas de trabajo futuro concretas y bien motivadas (no genéricas como ``aplicar \textit{deep learning}").}
\end{orientacion}

[Conclusiones aquí: 3--4 párrafos.]


% ============================================================
%  REFERENCIAS
% ============================================================

\begin{thebibliography}{99}

\bibitem{ref1}
A. Autor1 y B. Autor2,
``Título del artículo,''
\textit{Nombre de la Revista},
vol.~X, núm.~Y, pp.~ZZ--ZZ, a\~{n}o.
\href{https://doi.org/XX.XXXX/XXXXXXX}{DOI:XX.XXXX/XXXXXXX}

\bibitem{ref2}
C. Autor3,
``Título del artículo,''
in \textit{Proc. Nombre de la Conferencia (ABREV)},
Ciudad, País, a\~{n}o, pp.~ZZ--ZZ.

\bibitem{ref3}
D. Autor4, E. Autor5 y F. Autor6,
``Título del artículo,''
\textit{arXiv preprint arXiv:XXXX.XXXXX}, a\~{n}o.
\href{https://arxiv.org/abs/XXXX.XXXXX}{arXiv:XXXX.XXXXX}

\bibitem{ref4}
G. Autor7,
\textit{Título del libro},
N.\ ed.
Ciudad: Editorial, a\~{n}o.

% Añadir el resto de referencias aquí

\end{thebibliography}


% ============================================================
%  ANEXOS (opcionales)
% ============================================================

\appendix

\section{Contribución individual}

\begin{table}[H]
\centering
\caption{Contribución de cada miembro del grupo.}
\small
\begin{tabular}{ll}
\toprule
\textbf{Miembro} & \textbf{Contribución principal} \\
\midrule
{[}Nombre Apellido1{]} & {[}EDA, preprocesamiento, Sección~3{]} \\
{[}Nombre Apellido2{]} & {[}Implementación problemática 1, Sección~4.3{]} \\
{[}Nombre Apellido3{]} & {[}Implementación problemática 2, Sección~4.4{]} \\
\bottomrule
\end{tabular}
\end{table}

\section{Enlace al repositorio}

El código fuente del trabajo está disponible en:
\url{https://github.com/[usuario]/[nombre-repositorio]}

\section{Uso de herramientas de IA generativa}

\begin{orientacion}
\orientaciontitle
\textit{Indica si has utilizado herramientas de IA generativa (ChatGPT, GitHub Copilot, etc.) en la elaboración del trabajo, cómo las has usado y en qué partes. Si no se han utilizado, indícalo explícitamente.}
\end{orientacion}

[Descripción del uso, o bien: ``No se han utilizado herramientas de IA generativa en la elaboración de este trabajo".]

\end{document}